%% ======================================================================
%% ABSTRACT
%% ======================================================================
\begin{abstract}
Self-consistency samples several chains of thought and returns the majority
answer. This has been the traditional way to turn extra inference compute into accuracy on
reasoning and logic benchmarks, but it selects by answer frequency alone and is therefore
blind to whether the reasoning behind any candidate is valid. The accuracy declines sharply as the difficulty of the problems increases. For small open-weight models, the best answer that is already present among the four sampled chains of thought
beats their majority vote by six points on GSM8K, twelve on MATH-500 and eighteen on
JEEBench. Therefore, we introduce \textbf{\method{}}, \textbf{E}xplicit \textbf{S}trategy \textbf{V}erification, a gold-free procedure that never re-solves the problem: a zero-cost disagreement gate, one contrastive
verification call requiring each candidate's strategy to be bound to quoted givens and its preconditions checked, and a deterministic rule accepting only
candidates with no mismatched binding and no violated condition. Across 2,918 instances spanning GSM8K, MATH-500 and JEEBench, \method{} improves accuracy by 5.0, 5.9 and 3.5 points over CoT maj@4 (148 repairs against 4 regressions; $p \le 10^{-7}$ on all three), while Universal Self-Consistency reaches significance
on none. Substituting the solver for the stronger checker reverses the
effect on all three benchmarks, locating the method's requirement precisely: gold-free
verification needs a verifier stronger than the model under test. Every problem is paired with a
structural perturbation that either breaks the strategy the original invites or
preserves it, a design that manufactures selection errors of known cause and
separates strategy misfit from mere novelty---the error class frequency-based
selection cannot see.
\end{abstract}

%% ================================================
\section{Introduction}

Large language models solve multi-step reasoning problems by generating an
explicit chain of thought before committing to an answer~\cite{wei2022chain}. Because any single chain can go wrong, the standard
way to spend additional inference compute is to sample several of them and return their majority answer, a procedure known as self-consistency~\cite{wang2022self}. Yet majority voting selects among candidates by answer frequency alone, and is therefore blind to whether the reasoning that produced any of them is valid. The accuracy this forfeits is
substantial: for small open-weight models, the best answer already present among four sampled chains of thought exceeds their majority vote by six points on GSM8K, twelve points on MATH-500 and eighteen points on JEEBench. Existing gold-free remedies attack this gap in one of two ways, either re-solving the problem from a critique of the first attempt~\cite{madaan2023self,shinn2023reflexion}, or asking the model to nominate the most consistent of its own samples~\cite{chen2023universal}. The first pays full inference cost a second
time, and the second inherits the very biases that produced the error.

We argue that a practical gold-free selection method must satisfy three
requirements simultaneously: (1) it must never consult a reference answer, so
that it is deployable at inference time; (2) it must not re-solve the problem,
so that its cost is a fraction of the sampling budget it augments rather than a
multiple of it; and (3) it must not damage answers that were already correct,
since a method that fires on uncertainty can break a correct majority as easily
as repair a wrong one (Section~\ref{sec:requirements}). Existing methods satisfy
at most two. Re-solving approaches violate (2) by construction, and when the
model critiques its own work they violate (3) as well, because the blind spot
that produced the error also confirms it~\cite{huang2024large}. Selection
methods such as Universal Self-Consistency respect (1) and (2), but ask the
solver to adjudicate its own samples and so gain little over the majority vote
they replace. Satisfying all three at once is not a trivial issue.

To bridge this research gap, we propose \textbf{\method{}} (\textbf{Explicit Strategy
Verification}), a multi-stage gold-free pipeline that selects among answers the
solver has already produced and never re-solves the problem. A disagreement
gate first groups the samples into distinct candidates and halts immediately
when they agree, which costs no model calls and concentrates all verification
spend on the instances that can actually change
(Section~\ref{sec:gate}). Where they disagree, a single call presents the
distinct candidates---deduplicated, with vote counts hidden---to a checker
stronger than the solver, and a decision rule commits to one candidate or
abstains and keeps the majority answer (Sections~\ref{sec:verify}--\ref{sec:decide}). On
GSM8K, MATH-500 and JEEBench, across three open-weight solvers on the first and
two on the others, \method{} improves accuracy over majority voting by 5.0, 5.9
and 3.5 points respectively, against 0.9, 0.4 and $-1.9$ points for Universal
Self-Consistency, at 44\%, 83\% and 100\% of the sampling budget in additional
compute. To attribute the gain, we evaluate each stage in isolation on a
benchmark of structural perturbations built for the purpose.

To summarize, our contributions are:
\begin{itemize}
  \item We identify three requirements that a gold-free selection method must
        meet jointly, and show that no existing method meets all three.
  \item We propose \method{}, a gated verification pipeline that improves on
        majority voting and Universal Self-Consistency on every benchmark and
        solver we evaluate, without re-solving any problem.
  \item We construct a perturbation benchmark of matched problem pairs in which
        one variant breaks the strategy the original invites and the other
        preserves it, with an independently verified trap answer that makes the
        cause of each failure known.
  \item Component-wise experiments locate the gain in the checker's strength
        rather than in the verification format, and show that substituting the
        solver for the checker reverses the method's effect entirely.
\end{itemize}
